% =====================================================================
%  Trabalho de Redes -- IC/UFRJ
%  Bufferbloat, controle de congestionamento e a geometria da informacao do AIMD
%  Compila com pdflatex (2 passagens). No Overleaf compila sem ajustes.
% =====================================================================
\documentclass[11pt,a4paper]{article}

\usepackage[T1]{fontenc}
\usepackage[utf8]{inputenc}
\usepackage[brazilian]{babel}
\usepackage{lmodern}
\usepackage{amsmath,amssymb,mathtools}
\usepackage{microtype}
\usepackage[margin=2.5cm]{geometry}
\usepackage{booktabs}
\usepackage{tabularx}
\usepackage{xcolor}
\usepackage{enumitem}
\usepackage{tikz}
\usepackage{pgfplots}
\pgfplotsset{compat=1.18}
\usetikzlibrary{arrows.meta,positioning,calc,shapes.geometric}
\usepackage[most]{tcolorbox}
\usepackage{listings}
\usepackage[colorlinks=true,linkcolor=blue!50!black,urlcolor=blue!50!black,citecolor=blue!50!black]{hyperref}

% --------------------------------------------------------------------- parametros editaveis
\newcommand{\semestre}{2026/2}
\newcommand{\dataPreliminar}{[a definir]}   % versao preliminar + slides
\newcommand{\dataFinal}{[a definir]}        % versao final
\newcommand{\tamanhoGrupo}{3}
\newcommand{\linkCodigo}{https://drive.google.com/drive/folders/1OXwnsYDRsD-7qxpBDXE2uCf4oCdKQxyO?usp=sharing}

% --------------------------------------------------------------------- cores e marcadores de nivel
\definecolor{cfair}{RGB}{30,110,180}
\definecolor{ceff}{RGB}{200,80,30}
\definecolor{ctraj}{RGB}{40,40,40}
\definecolor{nv1}{RGB}{46,139,87}
\definecolor{nv2}{RGB}{30,110,180}
\definecolor{nv3}{RGB}{200,110,0}
\definecolor{nv4}{RGB}{160,30,60}

\newcommand{\tagnivel}[2]{\tikz[baseline=(t.base)]\node[fill=#1,text=white,rounded corners=2pt,inner sep=2.2pt,font=\footnotesize\bfseries\sffamily](t){#2};}
\DeclareRobustCommand{\basico}{\tagnivel{nv1}{BÁSICO}}
\DeclareRobustCommand{\inter}{\tagnivel{nv2}{INTERMEDIÁRIO}}
\DeclareRobustCommand{\avancado}{\tagnivel{nv3}{AVANÇADO}}
\DeclareRobustCommand{\pesquisa}{\tagnivel{nv4}{PESQUISA}}
\newcommand{\pontos}[1]{\noindent{\sffamily\small\color{gray!60!black}Valor: #1}\par\smallskip}

\newtcolorbox{caixa}[2][]{colback=#2!5,colframe=#2!70!black,fonttitle=\bfseries\sffamily,
  boxrule=0.6pt,arc=2pt,left=6pt,right=6pt,top=4pt,bottom=4pt,#1}

% --------------------------------------------------------------------- matematica
\newcommand{\R}{\mathbb{R}}
\newcommand{\AM}{\operatorname{AM}}
\newcommand{\GM}{\operatorname{GM}}
\newcommand{\Dis}{D_{\mathrm{IS}}}
\newcommand{\ones}{\mathbf{1}}

% --------------------------------------------------------------------- codigo
\lstdefinestyle{sh}{language=bash,basicstyle=\ttfamily\small,columns=fullflexible,
  keepspaces=true,backgroundcolor=\color{gray!8},frame=single,rulecolor=\color{gray!40},
  breaklines=true,showstringspaces=false,commentstyle=\color{gray!70!black},literate={-}{{-}}1}
\lstdefinestyle{py}{language=Python,basicstyle=\ttfamily\small,columns=fullflexible,
  keepspaces=true,backgroundcolor=\color{gray!8},frame=single,rulecolor=\color{gray!40},
  breaklines=true,showstringspaces=false,commentstyle=\color{gray!70!black},literate={-}{{-}}1,
  keywordstyle=\color{blue!60!black}\bfseries}

\setlist{itemsep=2pt,topsep=4pt}
\setlength{\parskip}{3pt}
\setlength{\emergencystretch}{3em}
\makeatletter\g@addto@macro\UrlBreaks{\do\-\do\.\do\_}\makeatother

\title{\bfseries Trabalho de Redes \semestre{} -- IC/UFRJ\\[4pt]
\large Bufferbloat, controle de congestionamento e\\ a geometria da informação do AIMD}
\author{Prof. Daniel Menasché}
\date{}

\begin{document}
\maketitle
\vspace{-1.2cm}

% =====================================================================
\begin{caixa}[title=Resumo do trabalho em uma caixa]{gray}
Você vai usar o Mininet para emular uma rede com um roteador gargalo e medir, no
\emph{kernel Linux real}, a janela de congestionamento de vários fluxos TCP. O trabalho
começa com um estudo curto de \emph{bufferbloat}, para aprender o ambiente, e depois
usa uma lente nova --- a \textbf{geometria da informação do AIMD} --- para medir
\emph{justiça} e \emph{eficiência} como duas componentes ortogonais de uma única
distância. Várias previsões dessa teoria foram verificadas apenas em simulação
numérica; \textbf{nunca foram testadas contra uma pilha TCP real}. É aí que vocês
entram.

\smallskip
\begin{tabularx}{\linewidth}{@{}lX@{}}
\basico & tarefas guiadas, com passo a passo; todos devem conseguir fazer.\\
\inter & exige escrever código de medição e análise; respostas não são únicas.\\
\avancado & exige desenho experimental próprio e interpretação cuidadosa.\\
\pesquisa & problemas em aberto; bons resultados podem entrar em artigo (Seção~\ref{sec:coautoria}).
\end{tabularx}
\end{caixa}

\section*{Entrega}
\begin{itemize}
\item Versão preliminar, com slides para serem apresentados: \dataPreliminar.
\item Versão final, levando em conta os comentários: \dataFinal.
\item Grupos de até \tamanhoGrupo{} alunos. Cada grupo mantém um repositório no GitHub
desde a primeira semana (Seção~\ref{sec:repro}).
\end{itemize}

\tableofcontents

% =====================================================================
\section{Visão geral}

Começamos com um estudo curto do fenômeno do \emph{bufferbloat} com TCP Reno: ele serve
para montar o ambiente e aprender a medir a janela de congestionamento, a fila do
roteador e o RTT, que são a matéria-prima de todo o resto. Essa parte segue o roteiro do projeto
de Kurtis Heimerl e Matt Johnson para a disciplina CSE~461 da Universidade de Washington
(\url{https://courses.cs.washington.edu/courses/cse461/22wi/projects/project3/}).

Em seguida, usamos o mesmo ambiente para fazer perguntas mais finas sobre a
\emph{dinâmica} do controle de congestionamento:
\begin{itemize}
\item Quando dois fluxos Reno começam com janelas diferentes, quão rápido eles ficam
“justos”? A teoria prevê que a injustiça cai por um fator $\beta^2$ a cada ciclo de
perda ($\beta=1/2$ no Reno). Isso vale no Linux?
\item Por que um fluxo com RTT duas vezes maior recebe às vezes metade da banda e às
vezes um quarto? A teoria diz que a resposta depende do \emph{grau de sincronização}
das perdas, com uma fórmula explícita. Ninguém mediu isso no Mininet.
\item O CUBIC tem expoente de justiça $4/3$ e o Reno tem $2$. O que isso significa, e
como se mede?
\end{itemize}

O ponto de partida teórico é o manuscrito \emph{The Information Geometry of AIMD}
(A.~Bhaya, D.~S.~Menasché e H.~Mirandola), fornecido junto com este enunciado,
com slides em português e dois scripts Python de verificação numérica. O manuscrito
está em preparação: \textbf{não o redistribua fora da turma}. A Seção~\ref{sec:primer}
resume tudo o que você precisa dele; não é preciso entender todas as provas para fazer
o trabalho.

\subsection*{Metas de aprendizado}
\begin{itemize}
\item Visualizar o “dente de serra” da janela de congestionamento e a ocupação de
buffer do roteador; entender por que buffers grandes levam a desempenho ruim.
\item Entender AIMD como \emph{dois movimentos geométricos}: o aumento aditivo melhora
a justiça, a redução multiplicativa mexe só na eficiência.
\item Medir justiça, eficiência, sincronização de perdas e expoentes de escala em
experimentos reais, com barras de erro.
\item Produzir experimentos \emph{reprodutíveis}: qualquer pessoa deve regenerar seus
gráficos com um único comando.
\end{itemize}

% =====================================================================
\section{Leituras}
\label{sec:leituras}

\paragraph{Bufferbloat e Mininet.}
\begin{itemize}
\item J.~Gettys e K.~Nichols, \emph{Bufferbloat: Dark Buffers in the Internet}, ACM
Queue, 2011. Ver também \url{https://pt.wikipedia.org/wiki/Bufferbloat}.
\item Mininet walkthrough: \url{http://mininet.org/walkthrough/} e a API Python:
\url{http://mininet.org/api/annotated.html}.
\item N.~Handigol et al., \emph{Reproducible network experiments using
container-based emulation}, ACM CoNEXT, 2012 (por que e quando o Mininet é confiável).
\end{itemize}

\paragraph{Controle de congestionamento.}
\begin{itemize}
\item Kurose \& Ross, \emph{Redes de Computadores e a Internet}, Seção 3.7 (TCP) e
Seção 4.2.4 (tamanho de buffer).
\item D.-M.~Chiu e R.~Jain, \emph{Analysis of the increase and decrease algorithms for
congestion avoidance}, 1989 --- origem do diagrama da Figura~\ref{fig:chiujain}.
\item M.~Mathis et al., \emph{The macroscopic behavior of the TCP congestion avoidance
algorithm}, 1997 --- a “lei da raiz quadrada”.
\item G.~Appenzeller, I.~Keslassy e N.~McKeown, \emph{Sizing router buffers}, ACM
SIGCOMM 2004, \url{https://dl.acm.org/doi/10.1145/1030194.1015499}.
\item S.~Ha, I.~Rhee e L.~Xu, \emph{CUBIC: a new TCP-friendly high-speed TCP variant},
2008; e a RFC~9438 (CUBIC).
\end{itemize}

\paragraph{Geometria da informação do AIMD.}
\begin{itemize}
\item O manuscrito (\texttt{main.tex}) e os slides (\texttt{slides-pt.tex}). Leia
primeiro os slides, depois as Seções 3, 6, 7 e 8 do manuscrito.
\item Para quem quiser ir além: S.~Amari, \emph{Information Geometry and Its
Applications}, Springer, 2016, Caps.~1--2.
\end{itemize}

% =====================================================================
\section{A teoria em quatro páginas}
\label{sec:primer}

Esta seção é um resumo autocontido. Os resultados citados estão provados no manuscrito;
aqui o objetivo é saber \emph{o que medir} e \emph{o que esperar}.

\subsection{A família binomial e o expoente de justiça}

Um protocolo da família binomial (Bansal--Balakrishnan) mantém uma janela $w$ e faz
\begin{equation}
\text{a cada RTT sem perda:}\quad w\leftarrow w+a\,w^{-k},
\qquad
\text{a cada perda:}\quad w\leftarrow w-b\,w^{l}.
\label{eq:binomial}
\end{equation}
O TCP Reno é $(k,l)=(0,1)$ com $b=1/2$: soma 1 pacote por RTT e divide a janela por
2 na perda. Escrevemos $\beta$ para a \textbf{fração retida} da janela na perda
($w\leftarrow\beta w$): $\beta=0{,}5$ no Reno e $\beta=0{,}7$ no CUBIC do Linux.

Se as perdas são independentes por pacote, com probabilidade $p$, o equilíbrio do
modelo fluido satisfaz $p\propto w^{-\alpha}$, ou seja,
\begin{equation}
\bar w\;\propto\;p^{-1/\alpha},\qquad \alpha=k+l+1 .
\label{eq:alpha}
\end{equation}
O número $\alpha$ é o \textbf{expoente de justiça}: o equilíbrio maximiza
$\sum_r U(x_r)$ com $U'(x)\propto x^{-\alpha}$ (justiça $\alpha$ de Mo--Walrand).
$\alpha=1$ é justiça proporcional, $\alpha=2$ é “atraso potencial mínimo”, e
$\alpha\to\infty$ é max-min.

\begin{center}
\begin{tabular}{@{}lcccl@{}}
\toprule
Protocolo & $k$ & $l$ & $\alpha$ & previsão $\bar w\propto p^{-1/\alpha}$\\
\midrule
Reno (AIMD) & $0$ & $1$ & $2$ & $p^{-1/2}$ (lei de Mathis)\\
Scalable TCP (MIMD) & $-1$ & $1$ & $1$ & $p^{-1}$\\
AIAD & $0$ & $0$ & $1$ & $p^{-1}$\\
CUBIC (correspondência de escala) & $-2/3$ & $1$ & $4/3$ & $p^{-3/4}$\\
\bottomrule
\end{tabular}
\end{center}

\paragraph{RTTs diferentes.} Com $a_r\propto1/T_r^2$ (o RTT entra duas vezes: na
taxa de incremento e na conversão janela$\to$taxa), a taxa de equilíbrio do fluxo $r$ é
\begin{equation}
x_r\;\propto\;T_r^{-2/\alpha_{\rm ef}} .
\label{eq:rttbias}
\end{equation}
Para o Reno com perdas \emph{dessincronizadas}, $\alpha_{\rm ef}=2$ e $x_r\propto1/T_r$
(viés linear em RTT). Com perdas \emph{sincronizadas}, $\alpha_{\rm ef}=1$ e
$x_r\propto1/T_r^2$ (viés quadrático). Veja a Seção~\ref{sec:sync}.

\subsection{AIMD são dois passos de tamanho fixo, em duas geometrias}

Na métrica euclidiana, o aumento aditivo é um passo de tamanho fixo $a$. Na métrica
invariante por escala $ds=dx/x$ (a única métrica separável invariante por dilatações,
que coincide com a métrica de Fisher de uma família de escala), a redução
multiplicativa $x\mapsto\beta x$ percorre sempre a mesma distância
\[
d(\beta x,x)=\int_{\beta x}^x\frac{du}{u}=|\log\beta|,
\]
\emph{independentemente de $x$}. A assimetria entre aumento e redução de Chiu--Jain é
a assimetria entre translações e dilatações.

\subsection{Justiça e eficiência como componentes ortogonais}

Considere $n$ fluxos com taxas $x=(x_1,\dots,x_n)$ dividindo um gargalo de capacidade
$c$, e seja $\mu$ a direção que o protocolo considera justa ($\mu=\ones$ quando todos
têm o mesmo RTT). O ótimo é $x^\star=c\,\mu/\sum_r\mu_r$. Defina $z_r=x_r/\mu_r$,
$\AM$ e $\GM$ as médias aritmética e geométrica, e a \textbf{divergência de
Itakura--Saito}
\begin{equation}
\Dis(x\|y)=\sum_{r=1}^n\Big[\frac{x_r}{y_r}-\log\frac{x_r}{y_r}-1\Big]\;\ge0 .
\end{equation}

\begin{caixa}[title=Teorema (decomposição pitagórica)]{cfair}
Com a \emph{utilização projetada} $\hat\rho=\AM(z)\sum_r\mu_r/c$,
\begin{equation}
\Dis(x\|x^\star)\;=\;
\underbrace{n\log\frac{\AM(z)}{\GM(z)}}_{\textstyle\mathcal U(x)\ \text{(injustiça)}}
\;+\;
\underbrace{n\big(\hat\rho-\log\hat\rho-1\big)}_{\textstyle\mathcal I(x)\ \text{(ineficiência)}} .
\label{eq:pitagoras}
\end{equation}
Os dois termos são $\ge0$ e se anulam juntos só em $x=x^\star$. $\mathcal U$ é
invariante por escala: não muda se todos os fluxos forem multiplicados pelo mesmo
fator.
\end{caixa}

Com $\mu=\ones$ a interpretação é imediata (Figura~\ref{fig:chiujain}): o ponto
$\hat x=\AM(x)\,\ones$ é a projeção de $x$ na reta justa \emph{mantendo a vazão total};
$\mathcal U$ mede o quanto $x$ está fora da reta justa, e $\mathcal I$ o quanto $\hat x$
está longe de $x^\star$ ao longo dela.

\begin{figure}[t]
\centering
\begin{tikzpicture}[scale=0.55]
  \draw[->,gray] (0,0) -- (10.8,0) node[right,black,font=\small] {$x_1$};
  \draw[->,gray] (0,0) -- (0,10.8) node[above,black,font=\small] {$x_2$};
  \draw[ceff,very thick] (0.4,9.6) -- (9.6,0.4)
      node[pos=0.04,above right,font=\scriptsize,black] {eficiente: $x_1+x_2=c$};
  \draw[cfair,very thick,dashed] (0,0) -- (10,10)
      node[pos=0.97,below right,font=\scriptsize,black] {justa: $x_1=x_2$};
  \draw[ctraj,thick,-{Stealth[length=2mm]}]
    (8.500,1.500) -- (4.250,0.750) -- (6.750,3.250) -- (3.375,1.625) --
    (5.875,4.125) -- (2.938,2.062) -- (5.438,4.562) -- (2.719,2.281) --
    (5.219,4.781) -- (2.609,2.391) -- (5.109,4.891);
  \fill[cfair] (5,5) circle (3pt) node[above right,font=\scriptsize,black] {$x^\star$};
  \fill[ctraj] (8.5,1.5) circle (2.4pt) node[right,font=\scriptsize,black] {$x_0$};
  % projecao de um ponto
  \coordinate (P) at (6.75,3.25);
  \coordinate (H) at (5,5);
  \node[font=\scriptsize,ceff] at (7.3,0.55) {MD: radial};
  \node[font=\scriptsize,cfair] at (2.5,4.2) {AI: direção $\ones$};
  % painel da direita: projecao
  \begin{scope}[xshift=14cm]
  \draw[->,gray] (0,0) -- (10.8,0) node[right,black,font=\small] {$x_1$};
  \draw[->,gray] (0,0) -- (0,10.8) node[above,black,font=\small] {$x_2$};
  \draw[ceff,very thick] (0.4,9.6) -- (9.6,0.4);
  \draw[cfair,very thick,dashed] (0,0) -- (10,10);
  \fill[ctraj] (6,1.5) circle (3pt) node[right,font=\scriptsize] {$x$};
  \fill[cfair!70!black] (3.75,3.75) circle (3pt) node[above left,font=\scriptsize] {$\hat x=\AM(x)\ones$};
  \fill[cfair] (5,5) circle (3pt) node[above right,font=\scriptsize] {$x^\star$};
  \draw[gray,densely dotted,thick] (6,1.5) -- (3.75,3.75)
     node[midway,below left,font=\scriptsize,black] {$\mathcal U$};
  \draw[ceff,line width=1.6pt] (3.75,3.75) -- (5,5)
     node[midway,above left,font=\scriptsize,black] {$\mathcal I$};
  \end{scope}
\end{tikzpicture}
\caption{Esquerda: trajetória do AIMD sincronizado para dois fluxos (diagrama de
Chiu--Jain). A redução multiplicativa (MD) é radial: preserva as proporções e,
portanto, \emph{não muda} $\mathcal U$. O aumento aditivo (AI) anda na direção $\ones$:
reduz $\mathcal U$. Direita: a decomposição \eqref{eq:pitagoras}. A projeção $\hat x$
tem a mesma vazão total que $x$; $\mathcal U$ é a parte “transversal” e $\mathcal I$ a
parte “ao longo da reta justa”. Os comprimentos no desenho são ilustrativos: as
distâncias verdadeiras são divergências, não distâncias euclidianas.}
\label{fig:chiujain}
\end{figure}

\paragraph{O que acontece em cada ciclo.} Para o Reno com perdas sincronizadas, amostrando
o estado logo antes de cada perda:
\begin{enumerate}[label=(\roman*)]
\item a MD $x\mapsto\beta x$ deixa $\mathcal U$ \emph{exatamente} igual e leva
$\hat\rho\mapsto\beta\hat\rho$;
\item a AI reduz $\mathcal U$ estritamente;
\item de um ciclo para o outro, $\mathcal U_{m+1}/\mathcal U_m\to\beta^2$
($=1/4$ no Reno);
\item no regime permanente, $\mathcal U\to0$ e $\mathcal I$ oscila entre $0$ e
$n(\beta-\log\beta-1)$ --- é a sonda de banda, que \emph{deve} persistir.
\end{enumerate}

\subsection{Sincronização de perdas}
\label{sec:sync}

Modele cada evento de congestionamento assim: com probabilidade $s$ ele é
\emph{global} (todos os $n$ fluxos reduzem a janela); com probabilidade $1-s$ ele é
\emph{individual} (um único fluxo reduz, escolhido com probabilidade proporcional à sua
taxa). O manuscrito mostra que o expoente de justiça efetivo do Reno é
\begin{equation}
\alpha_{\rm ef}(s)\;=\;1+\gamma(s),\qquad
\gamma(s)=\frac{1-s}{s\,n+1-s},
\label{eq:interp}
\end{equation}
que vai de $\alpha=2$ ($s=0$, dessincronizado) a $\alpha=1$ ($s=1$, sincronizado). O
ponto médio está em $s=1/n$, e não em $s=1/2$: com muitos fluxos, \emph{uma pequena
fração de eventos globais já domina}.

\paragraph{Como estimar $s$ a partir de um traço.} Agrupe as reduções de janela de todos
os fluxos em \emph{eventos} (reduções que acontecem dentro de uma janela de tempo
$\Delta$, da ordem de um RTT). Se $\bar m$ é o número médio de fluxos atingidos por
evento, o modelo dá $\bar m=s\,n+(1-s)$, logo
\begin{equation}
\hat s=\frac{\bar m-1}{n-1}.
\label{eq:shat}
\end{equation}

\subsection{O que a teoria \emph{não} cobre}
Tudo acima é para o modelo fluido \emph{sem atraso de realimentação}, com perdas
idealizadas. Ficam de fora: (a) estabilidade com atraso; (b) a dinâmica completa do
CUBIC, que tem um estado extra $W_{\max}$; (c) o dente de serra em redes com mais de um
gargalo; (d) detalhes do Linux como \emph{slow start},
HyStart, PRR, SACK, pacing e TSO. Medir o quanto cada um desses itens desvia das
previsões é exatamente o tipo de resultado que interessa a um artigo.

% =====================================================================
\section{Parte 1: Configuração \texorpdfstring{\basico}{}}
\label{sec:setup}

\paragraph{VM.} Siga \url{https://github.com/kaichengyan/mininet-vagrant} para montar
a VM Vagrant (Ubuntu 20.04 com Mininet). Alternativamente, uma VM Ubuntu 22.04/24.04 com
\texttt{sudo apt install mininet} também serve; registre qual você usou
(Seção~\ref{sec:repro}).

\paragraph{Código inicial.} No diretório inicial da VM:
\begin{lstlisting}[style=sh]
cd ~
wget --no-check-certificate 'https://docs.google.com/uc?export=download&id=12A0GGl7EJvKhBogz6H_fbzH33ZQT9FBQ' -O bufferbloat.zip
unzip bufferbloat.zip
\end{lstlisting}
O código também está em \url{\linkCodigo}. Os arquivos são:

\begin{center}
\begin{tabularx}{\textwidth}{@{}lX@{}}
\toprule
Arquivo & Propósito\\
\midrule
\texttt{run.sh} & Executa o experimento e gera os gráficos de uma vez.\\
\texttt{bufferbloat.py} & Cria a topologia e mede cwnd, fila e RTT.\\
\texttt{monitor.py} & Monitora o comprimento da fila.\\
\texttt{plot\_queue.py}, \texttt{plot\_ping.py} & Plotam a fila e o RTT.\\
\texttt{plot\_defaults.py}, \texttt{helper.py} & Utilitários para gráficos.\\
\texttt{webserver.py}, \texttt{index.html} & Servidor web (não usado nesta versão).\\
\bottomrule
\end{tabularx}
\end{center}

\paragraph{Pacotes e verificações.} Na VM:
\begin{lstlisting}[style=sh]
sudo apt-get update
sudo apt install python3-pip iperf3 ethtool
sudo python3 -m pip install mininet matplotlib numpy pandas
# algoritmos de controle de congestionamento disponiveis
sysctl net.ipv4.tcp_available_congestion_control
sudo modprobe tcp_scalable
# nao reaproveitar ssthresh/cwnd de conexoes anteriores entre rodadas
sudo sysctl -w net.ipv4.tcp_no_metrics_save=1
\end{lstlisting}
Rode também os scripts do manuscrito (\texttt{verify\_aimd.py} e \texttt{verify\_cubic.py}, com \texttt{python3}); eles só dependem de \texttt{numpy}.

\paragraph{Código da segunda metade.} A partir da Parte 2 você também usa o pacote
\texttt{ig\_starter} (\url{https://github.com/danielsadoc/trabalho-redes-ig}), descrito no Apêndice~\ref{app:codigo}:
\texttt{dumbbell.py} monta a topologia das Partes 4--6 e roda o experimento,
\texttt{cwnd\_monitor.py} amostra a janela com \texttt{ss}, \texttt{ig\_metrics.py} tem as
métricas da Seção~\ref{sec:primer} e \texttt{plot\_cwnd.py} faz o gráfico básico.

\paragraph{Observações.} Execute sempre com \texttt{sudo ./run.sh}. Se o Mininet
terminar com erro (ou com Control-C), rode \texttt{sudo mn -c} antes de recomeçar.

% =====================================================================
\section{Parte 2: Bufferbloat e o dente de serra do Reno \texorpdfstring{\basico}{}}
\pontos{1,5 ponto}
\label{sec:bufferbloat}

Esta parte é o pré-requisito prático de todo o resto: montar uma topologia no Mininet e
medir, ao mesmo tempo, a janela de congestionamento, a fila do roteador e o RTT.

\begin{figure}[htbp]
\centering
\begin{tikzpicture}[node distance=2.6cm,
  host/.style={draw,rounded corners,minimum width=1.3cm,minimum height=0.8cm,font=\small,fill=gray!10},
  sw/.style={draw,circle,minimum size=0.9cm,font=\small,fill=cfair!15}]
  \node[host] (h1) {h1};
  \node[sw,right=of h1] (s0) {s0};
  \node[host,right=of s0] (h2) {h2};
  \draw[thick] (h1) -- node[above,font=\scriptsize]{1 Gb/s} (s0);
  \draw[thick,ceff] (s0) -- node[above,font=\scriptsize]{1,5 Mb/s (gargalo)}
     node[below,font=\scriptsize]{buffer $q$ pacotes} (h2);
\end{tikzpicture}
\caption{Topologia da Parte 2. RTT mínimo entre h1 e h2: 20\,ms.}
\label{fig:topo1}
\end{figure}

Crie a topologia da Figura~\ref{fig:topo1} na classe \texttt{BBTopo} de
\texttt{bufferbloat.py}: h1 tem enlace rápido (1\,Gb/s) até o roteador, que tem uplink
lento (1,5\,Mb/s). O RTT de propagação entre h1 e h2 é 20\,ms e o buffer do roteador tem
$q=100$ pacotes (150\,kB com MTU de 1500 bytes). Simultaneamente:
\begin{itemize}
\item inicie um fluxo TCP Reno de longa duração de h1 para h2 (iperf3);
\item colete 10 pings por segundo de h1 para h2 e registre os RTTs;
\item registre a ocupação da fila do gargalo (\texttt{monitor.py}) e a janela de
congestionamento do fluxo (\texttt{cwnd\_monitor.py}, do código da segunda metade).
\end{itemize}
Plote as séries temporais do RTT, da fila e da cwnd. Repita com $q=20$.

\paragraph{Gráficos obrigatórios} (6 no total):
\begin{itemize}
\item \path{reno-buffer-q100.png}, \path{reno-rtt-q100.png}, \path{reno-cwnd-q100.png};
\item \path{reno-buffer-q20.png}, \path{reno-rtt-q20.png}, \path{reno-cwnd-q20.png}.
\end{itemize}

\subsection*{Perguntas da Parte 2}
\begin{enumerate}[label=2.\arabic*]
\item Qual o RTT médio do ping e seu desvio padrão para $q=20$ e $q=100$? Escreva a
relação entre RTT e ocupação da fila (dica: atraso de fila $=$ ocupação / taxa de
serviço) e confira com os seus dados.
\item No gráfico de cwnd, meça a janela no pico e no vale do dente de serra. Compare o
pico com $\mathrm{BDP}+q$ e a razão vale/pico com $\beta=1/2$. Você vai usar
$c=\mathrm{BDP}+B$ como “capacidade” na Parte 4.
\item Por que um buffer grande aumenta o atraso sem aumentar a vazão? Qual o menor
buffer que mantém o gargalo sempre ocupado com um único fluxo Reno (regra
$B=\mathrm{RTT}\cdot C$)? Confira nos seus dados.
\item O livro (Seção 4.2.4) cita a regra $B=\mathrm{RTT}\cdot C/\sqrt N$ de Appenzeller
et al. Ela supõe que os $N$ fluxos estão \emph{dessincronizados}. Por quê? O que
aconteceria com a regra se as perdas fossem perfeitamente sincronizadas? (Você vai
medir isso na Parte 6.)
\end{enumerate}

% =====================================================================
\section{Parte 3: Aquecimento teórico e computacional \texorpdfstring{\basico}{}\ \texorpdfstring{\inter}{}}
\pontos{1,5 ponto}
\label{sec:aquecimento}

Estes exercícios garantem que vocês sabem calcular as quantidades que serão medidas.
Respostas curtas, com as contas.

\begin{enumerate}[label=3.\arabic*]
\item \basico\ Usando $\alpha=k+l+1$, preencha a tabela da Seção~\ref{sec:primer} para
AIAD, MIMD e para o binomial $(k,l)=(1,1)$ (IIAD). Qual deles é o mais “justo” no
sentido de $\alpha$?
\item \basico\ Para $n=2$ e $\mu=\ones$, mostre que $\mathcal U(x)=2\log\frac{(x_1+x_2)/2}{\sqrt{x_1x_2}}$
e que $\mathcal U(\lambda x)=\mathcal U(x)$ para todo $\lambda>0$. Conclua que a MD
sincronizada não muda $\mathcal U$.
\item \basico\ Calcule a amplitude $n(\beta-\log\beta-1)$ da oscilação de $\mathcal I$
no regime permanente para Reno e CUBIC com $n=2$ e $n=3$.
\item \inter\ Para $n=2$, $\mu=\ones$ e $x=(1+u,1-u)$ com $u$ pequeno, mostre que
$\mathcal U\approx u^2$. Use isso e o mapa $u\mapsto\beta u$ para explicar por que a
razão por ciclo é $\beta^2$ e não $\beta$.
\item \inter\ Leia as funções \texttt{unfairness}, \texttt{inefficiency} e
\texttt{itakura\_saito} de \texttt{ig\_metrics.py} e confira que implementam as fórmulas
da Seção~\ref{sec:primer}. Escreva um teste (\texttt{test\_ig\_metrics.py}) que verifica a
identidade \eqref{eq:pitagoras} em 1000 pontos aleatórios com erro $<10^{-10}$, para
$\mu=\ones$ e para $\mu=(1,1/4,1/9)$, e que verifica $\mathcal U(\lambda x)=\mathcal U(x)$.
\item \inter\ Rode \texttt{verify\_aimd.py}. Explique, com suas palavras, o que o item 3
da saída verifica e por que a razão por ciclo começa em $\approx0{,}33$ e converge para
$0{,}25$.
\end{enumerate}

% =====================================================================
\section{Parte 4: Medindo a decomposição pitagórica \texorpdfstring{\inter}{}}
\pontos{2,5 pontos}
\label{sec:pitagoras}

\begin{figure}[htbp]
\centering
\begin{tikzpicture}[
  host/.style={draw,rounded corners,minimum width=1.1cm,minimum height=0.7cm,font=\small,fill=gray!10},
  sw/.style={draw,circle,minimum size=0.9cm,font=\small,fill=cfair!15}]
  \node[host] (a1) at (0,1.6) {h1};
  \node[host] (a2) at (0,0.5) {h2};
  \node[font=\small] at (0,-0.35) {$\vdots$};
  \node[host] (an) at (0,-1.2) {h$n$};
  \node[sw] (s1) at (3.2,0.2) {s1};
  \node[sw] (s2) at (7.6,0.2) {s2};
  \node[host] (rx) at (10.2,0.2) {r};
  \draw[thick] (a1) -- node[above,sloped,font=\scriptsize]{$d_1$} (s1);
  \draw[thick] (a2) -- node[above,sloped,font=\scriptsize]{$d_2$} (s1);
  \draw[thick] (an) -- node[below,sloped,font=\scriptsize]{$d_n$} (s1);
  \draw[thick,ceff] (s1) -- node[above,font=\scriptsize]{gargalo $C$}
     node[below,font=\scriptsize]{buffer $B$ pacotes} (s2);
  \draw[thick] (s2) -- (rx);
\end{tikzpicture}
\caption{Topologia “haltere” (\emph{dumbbell}) usada nas Partes 4 e 6. Na Parte 4 todos
os atrasos de acesso são iguais; na Parte 6 eles diferem.}
\label{fig:dumbbell}
\end{figure}

\paragraph{Montagem.} Use a topologia da Figura~\ref{fig:dumbbell} com $n$ emissores
h1,\dots,h$n$ e um receptor r. Valores de referência: $C=10$\,Mb/s, RTT de propagação
de 40\,ms \emph{igual} para todos os fluxos, $B\in\{20,100\}$ pacotes no buffer do
enlace s1$\to$s2 (só nele). Cada emissor abre \emph{um} fluxo iperf3 para r, com o
algoritmo escolhido por fluxo (\texttt{iperf3 -C reno}); o \texttt{dumbbell.py} já faz
isso. Os fluxos começam
\emph{escalonados}: o fluxo $i$ começa em $t=15(i-1)$\,s, para criar uma injustiça
inicial grande. Duração mínima: 120\,s depois do início do último fluxo.

\paragraph{Medição.} Amostre a cada 10\,ms, para cada fluxo, \texttt{cwnd},
\texttt{ssthresh} e \texttt{rtt} (o \texttt{dumbbell.py} chama o
\texttt{cwnd\_monitor.py} em cada emissor e grava um CSV por fluxo).
Como os RTTs de propagação são iguais, use $x_r=\texttt{cwnd}_r$ (em pacotes) e
$\mu=\ones$. Para a capacidade, use $c=\mathrm{BDP}+B$ em pacotes, e compare com a média
de $\sum_r\texttt{cwnd}_r$ imediatamente antes das perdas.

\paragraph{O que fazer.}
\begin{enumerate}[label=4.\arabic*]
\item \basico\ Plote $\texttt{cwnd}_r(t)$ de todos os fluxos no mesmo gráfico, para
$n=2$, Reno, $B=100$. Identifique visualmente os eventos de perda.
\item \inter\ Plote $\mathcal U(t)$, $\mathcal I(t)$ e $\Dis(x(t)\|x^\star)$ no mesmo
gráfico (escala log no eixo $y$ para $\mathcal U$). Confira numericamente a identidade
\eqref{eq:pitagoras} nos seus dados.
\item \inter\ Implemente \texttt{detect\_backoffs} em \texttt{ig\_metrics.py} para
detectar automaticamente os instantes de perda. Use as mudanças de
\texttt{ssthresh}: no Linux a cwnd cai aos poucos durante a recuperação (PRR,
RFC~6937), ao longo de cerca de um RTT, então quedas grandes entre duas amostras
seguidas são raras. Para cada perda, compare
$\mathcal U$ imediatamente antes e imediatamente depois. A previsão (i) da
Seção~\ref{sec:primer} --- MD não muda $\mathcal U$ --- vale? Quando falha, quantos
fluxos reduziram a janela naquele evento?
\item \inter\ Amostrando antes de cada perda, estime a razão
$\mathcal U_{m+1}/\mathcal U_m$ e compare com $\beta^2=0{,}25$. Use só os ciclos depois que
todos os fluxos saíram do \emph{slow start} e antes de $\mathcal U$ cair no nível do
ruído de medição (como a cwnd é inteira, para $n=2$ esse piso é da ordem de
$(2\bar w)^{-2}$). Quantos ciclos limpos você conseguiu?
\item \inter\ Repita para CUBIC. Se o CUBIC fosse um AIMD com $\beta=0{,}7$, a razão
seria $0{,}49$. O que você mede? Por que o CUBIC não precisa seguir essa previsão?
\item \inter\ Repita com $n=3$ e $n=4$, e com $B=20$. O que acontece com a fração de
perdas que atingem todos os fluxos ao mesmo tempo? Como isso afeta as previsões (i) e
(iii)?
\end{enumerate}

\paragraph{Entregáveis obrigatórios:}
\begin{itemize}
\item \path{p4-cwnd-reno-n2-B100.png}, \path{p4-decomp-reno-n2-B100.png} e \path{p4-decomp-cubic-n2-B100.png};
\item uma tabela com a razão média por ciclo (com IC 95\%) para cada combinação
(algoritmo, $n$, $B$) testada.
\end{itemize}

% =====================================================================
\section{Parte 5: Leis de escala janela \texorpdfstring{$\times$}{x} perda \texorpdfstring{\inter}{}}
\pontos{1,5 ponto}
\label{sec:escala}

O objetivo é medir o expoente $\alpha$ de \eqref{eq:alpha} no Linux.

\paragraph{Montagem.} Um único fluxo através do gargalo (\texttt{dumbbell.py --n 1 --loss}). Introduza
perda aleatória independente por pacote com probabilidade $p$ \emph{na interface de saída
do switch} (não no emissor: perdas na fila local do emissor são vistas pelo TCP de outra
forma). No TCLink, o parâmetro \texttt{loss} é em \emph{porcento}. Use uma capacidade alta
o bastante (por exemplo 50\,Mb/s) e um buffer grande o bastante para que quase todas as
perdas sejam as aleatórias. Valores de referência: RTT de 100\,ms e $p$ em pelo menos
sete valores log-espaçados entre $2\times10^{-4}$ e $2\times10^{-2}$.

\paragraph{O que fazer.}
\begin{enumerate}[label=5.\arabic*]
\item \basico\ Para Reno, meça a janela média $\bar w(p)$ e plote $\log\bar w$ contra
$\log p$. Ajuste uma reta e estime $\hat\alpha=-1/\text{inclinação}$, com intervalo de
confiança (repita cada ponto pelo menos 5 vezes). Compare com a lei de Mathis
$\bar w\approx1{,}22/\sqrt p$.
\item \inter\ Repita para Scalable TCP (\texttt{-C scalable}). A previsão é
$\hat\alpha\approx1$. Leia o código-fonte de \texttt{tcp\_scalable.c} do seu kernel e
identifique o incremento e a fração retida $\beta$.
\item \inter\ Repita para CUBIC. O CUBIC do Linux tem uma região “amigável ao TCP” em
que se comporta como um AIMD; fora dela, a previsão é $\hat\alpha=4/3$. Você deve ver
\emph{dois regimes}. Encontre experimentalmente o valor de $p$ (e de RTT) onde ocorre a
transição e compare com a fórmula da RFC~9438.
\item \inter\ Quantos eventos de perda cada ponto contém? Discuta o compromisso entre
duração do experimento e precisão de $\hat\alpha$. (O \texttt{verify\_aimd.py} usa
$2\times10^5$ RTTs por ponto; qual a duração equivalente no seu experimento?)
\end{enumerate}

\paragraph{Gráfico obrigatório:} \path{p5-scaling.png}, com os três algoritmos no
mesmo gráfico log-log, as retas teóricas e as barras de erro.

% =====================================================================
\section{Parte 6: Sincronização e justiça entre RTTs \texorpdfstring{\avancado}{}}
\pontos{2,0 pontos}
\label{sec:sincronizacao}

Esta é a parte com maior potencial de pesquisa: \textbf{a fórmula \eqref{eq:interp}
nunca foi testada em uma pilha TCP real.}

\paragraph{Montagem.} Topologia da Figura~\ref{fig:dumbbell} com $n=3$ fluxos Reno e
atrasos de acesso escolhidos para que os RTTs de propagação fiquem na razão $1:2:3$ (por
exemplo, 20, 40 e 60\,ms). Meça os RTTs com \texttt{ping} antes de ligar os fluxos.
Atenção: o RTT \emph{efetivo} inclui o atraso de fila, que é comum a todos os fluxos e
comprime as razões. Use o \texttt{srtt} médio medido por fluxo, não o de propagação.

\paragraph{O que fazer.}
\begin{enumerate}[label=6.\arabic*]
\item \inter\ Implemente \texttt{sync\_estimate} em \texttt{ig\_metrics.py}, o estimador
$\hat s$ de \eqref{eq:shat}: detecte as reduções de
janela de cada fluxo, agrupe-as em eventos com uma janela $\Delta$ e calcule $\bar m$.
Mostre como $\hat s$ depende de $\Delta$ (teste $\Delta$ de 0,5 a 3 RTTs) e justifique sua
escolha.
\item \avancado\ Estime $\hat\alpha_{\rm ef}$ a partir de \eqref{eq:rttbias}: ajuste
$\log x_r$ contra $\log\bar T_r$; a inclinação é $-2/\hat\alpha_{\rm ef}$.
\item \avancado\ Faça $\hat s$ variar e construa o gráfico de $\hat\alpha_{\rm ef}$
contra $\hat s$, com a curva teórica \eqref{eq:interp} sobreposta. Maneiras de mexer na
sincronização:
\begin{itemize}
\item tamanho do buffer $B$ (de $0{,}25$ a $4$ BDPs);
\item disciplina de fila: drop-tail (padrão) $\times$ RED (\texttt{dumbbell.py --red});
\item número de fluxos $n\in\{3,5,10\}$ (com RTTs log-espaçados);
\item uma pequena perda aleatória adicional (que cria eventos individuais);
\item tráfego de fundo com fluxos curtos (por exemplo, downloads periódicos com \texttt{curl}).
\end{itemize}
\item \avancado\ Para $n=10$, a teoria coloca o ponto médio ($\alpha_{\rm ef}=1{,}5$) em
$s=1/n=0{,}1$; para $n=30$, $s=0{,}1$ já dá $\alpha_{\rm ef}\approx1{,}23$, perto do
regime sincronizado. Seus dados confirmam que o ponto médio se desloca para $s=1/n$
quando $n$ cresce?
\item \avancado\ Volte à pergunta 2.4. Para quais valores de $B$ as perdas foram
efetivamente dessincronizadas no seu experimento? O que isso diz sobre a regra
$\mathrm{RTT}\cdot C/\sqrt N$?
\end{enumerate}

\paragraph{Gráfico obrigatório:} \path{p6-alpha-vs-s.png}: $\hat\alpha_{\rm ef}$
contra $\hat s$, com barras de erro nos dois eixos e a curva teórica para cada $n$
usado. Este é o gráfico que mais interessa para o artigo.

% =====================================================================
\section{Parte 7: Dois gargalos \texorpdfstring{\avancado}{}}
\pontos{1,0 ponto}
\label{sec:doisgargalos}

O manuscrito mostra que a identificação “a redução define a geometria; $(k,l)$ define a
utilidade $\alpha$-justa” e a função de Lyapunov $\Dis(x^\star\|x)$ valem em qualquer
topologia, não só com um gargalo. Esta parte testa isso na rede da
Figura~\ref{fig:parking}, com os atrasos de acesso ajustados para que os três fluxos
tenham o mesmo RTT de propagação.

\begin{figure}[htbp]
\centering
\begin{tikzpicture}[
  sw/.style={draw,circle,minimum size=0.9cm,font=\small,fill=cfair!15}]
  \node[sw] (A) at (0,0) {A};
  \node[sw] (B) at (4,0) {B};
  \node[sw] (Cn) at (8,0) {C};
  \draw[thick,ceff] (A) -- node[above,font=\scriptsize]{$c_1=10$ Mb/s} (B);
  \draw[thick,ceff] (B) -- node[above,font=\scriptsize]{$c_2=5$ Mb/s} (Cn);
  \draw[cfair,very thick,-{Stealth}] (-0.4,-0.75) -- node[below,font=\scriptsize]{fluxo 1} (8.4,-0.75);
  \draw[ctraj,very thick,-{Stealth}] (-0.4,0.75) -- node[above,font=\scriptsize]{fluxo 2} (4.2,0.75);
  \draw[ctraj,very thick,-{Stealth}] (3.8,1.35) -- node[above,font=\scriptsize]{fluxo 3} (8.4,1.35);
\end{tikzpicture}
\caption{Rede “parking lot” com dois gargalos e três fluxos (exemplo de
Uchida--Kurose usado no manuscrito).}
\label{fig:parking}
\end{figure}

\begin{enumerate}[label=7.\arabic*]
\item \inter\ Calcule a alocação ótima para $\alpha=1$ (justiça proporcional),
$\alpha=2$ e $\alpha\to\infty$ (max-min). Para conferir: com $\alpha=2$ e o modelo de
preços do manuscrito, $x^\star\approx(2{,}43;\ 7{,}57;\ 2{,}57)$\,Mb/s.
\item \avancado\ Meça a alocação com Reno. Qual $\alpha$ descreve melhor o resultado?
Repita com RTTs desiguais (sem compensar o caminho mais longo do fluxo 1).
\item \avancado\ Calcule $\Dis(x^\star\|x(t))$ ao longo do experimento. Ela é
monótona? Lembre que o manuscrito prova a monotonicidade só para o modelo fluido sem
atraso.
\end{enumerate}

% =====================================================================
\section{Parte 8: Trilha de pesquisa \texorpdfstring{\pesquisa}{}}
\pontos{bônus de até 2,0 pontos}
\label{sec:pesquisa}

Os itens abaixo são problemas em aberto. Não há resposta conhecida; um resultado
negativo bem feito (“a previsão falha, e eis por quê”) também é um resultado. Escolha
\emph{no máximo dois} e converse com o professor antes de investir muito tempo.

\begin{enumerate}[label=R\arabic*.]
\item \textbf{Validação da interpolação.} Transformar o gráfico da Parte 6 num resultado
de artigo: vários $n$, repetições suficientes, análise de sensibilidade de $\Delta$, e
uma explicação para os desvios. Primeiro passo: 10 repetições de cada ponto da Parte 6.
\item \textbf{A contração transversal do CUBIC.} O CUBIC tem estado bidimensional
$(w,W_{\max})$ e o manuscrito deixa em aberto sua geometria. Meça
$\mathcal U_{m+1}/\mathcal U_m$ do CUBIC em função de $\beta$ e $C$ (parâmetros do módulo
\texttt{tcp\_cubic}) e proponha uma fórmula. Compare com uma simulação do mapa de ciclos
usando \texttt{verify\_cubic.py} como base.
\item \textbf{Atraso de realimentação.} Aumente o RTT (até 400\,ms) mantendo o BDP fixo
e verifique se $\Dis(x^\star\|x(t))$ deixa de ser monótona. Existe um RTT a partir do
qual a monotonicidade quebra? Isso dialoga com a questão aberta 1 do manuscrito.
\item \textbf{Disciplinas de fila com $\gamma>1$.} O manuscrito pergunta se existem
disciplinas de fila que realizam $\alpha>k+l+1$. Proponha uma (por exemplo, descarte com
probabilidade que cresce mais que linearmente na taxa do fluxo), implemente-a num
simulador simples em Python ou como qdisc, e meça $\alpha_{\rm ef}$.
\item \textbf{Dente de serra com dois gargalos.} Na rede da Figura~\ref{fig:parking},
procure empiricamente uma quantidade transversal que seja monótona de ciclo a ciclo
(questão aberta 5 do manuscrito). Candidata natural: $\mathcal U$ com $\mu=x^\star$.
\end{enumerate}

% =====================================================================
\section{Reprodutibilidade e formato dos dados}
\label{sec:repro}

Um dos objetivos do trabalho é construir um sistema fácil de rodar. Como os dados podem
entrar em um artigo, as regras abaixo são \textbf{obrigatórias} a partir da Parte 4.

\begin{itemize}
\item \textbf{Um comando por parte}: \texttt{sudo ./run.sh}, \texttt{sudo ./run\_p4.sh},
\texttt{sudo ./run\_p5.sh}, \dots, e \texttt{./plot\_all.sh} para regenerar todos os
gráficos a partir dos dados brutos, sem rodar o Mininet de novo.
\item \textbf{Dados brutos em CSV}, no formato gravado pelo \texttt{cwnd\_monitor.py}, com colunas
\texttt{t\_s, flow\_id, cc, cwnd\_pkts, ssthresh\_pkts, srtt\_ms, bytes\_acked} e um
\texttt{metadata.json} ao lado (o \texttt{dumbbell.py} já grava um) com: versão do kernel (\texttt{uname -r}), versão do
Mininet e do iperf3, topologia, $C$, $B$, atrasos, qdisc, algoritmo por fluxo, semente,
data e hora, e o \emph{commit} do repositório.
\item \textbf{Repetições}: pelo menos 5 por configuração, com intervalos de confiança de
95\%. Diga explicitamente quais rodadas foram descartadas e por quê.
\item \textbf{Sanidade}: em cada rodada, verifique que a CPU da VM não saturou
(\texttt{top} ou \texttt{mpstat}) e que a vazão total bate com $C$. O Mininet só é fiel
enquanto a CPU não é o gargalo.
\item \textbf{Repositório GitHub} desde a primeira semana, com \texttt{README.md}
explicando como reproduzir cada gráfico, e uma licença para o código (por exemplo, MIT)
e para os dados (por exemplo, CC-BY 4.0).
\end{itemize}

% =====================================================================
\section{O que entregar}
\label{sec:entrega}

\begin{itemize}
\item \textbf{Código}: todos os arquivos fornecidos e os seus, num único ZIP, e o link do
repositório GitHub.
\item \textbf{Relatório em PDF} com instruções para reproduzir os resultados e as
respostas às perguntas, identificadas pelo número (2.1, 3.4, 6.3, \dots).
\item \textbf{Gráficos} com os nomes pedidos em cada parte, também enviados como
arquivos separados no Google Classroom.
\item \textbf{Slides} para a apresentação da versão preliminar.
\item \textbf{Declaração de contribuição}: um parágrafo dizendo o que cada membro do
grupo fez.
\end{itemize}
Submeta via Google Classroom: PDF e imagens como arquivos isolados, e o código num ZIP
à parte.

\subsection*{Pontuação}
\begin{center}
\begin{tabular}{@{}llr@{}}
\toprule
Parte & Nível & Pontos\\
\midrule
1. Configuração & \basico & pré-requisito\\
2. Bufferbloat e dente de serra & \basico & 1,5\\
3. Aquecimento & \basico\ \inter & 1,5\\
4. Decomposição pitagórica & \inter & 2,5\\
5. Leis de escala & \inter & 1,5\\
6. Sincronização e RTTs & \avancado & 2,0\\
7. Dois gargalos & \avancado & 1,0\\
\midrule
8. Trilha de pesquisa & \pesquisa & bônus até 2,0\\
\bottomrule
\end{tabular}
\end{center}
A nota final é limitada a 10. Em cada parte, metade da nota vem dos experimentos e
gráficos, e metade das respostas e da interpretação. Uma resposta que explica
honestamente por que um experimento não deu certo vale mais do que um gráfico bonito
sem explicação.

% =====================================================================
\section{Coautoria e uso dos resultados}
\label{sec:coautoria}

Os experimentos deste trabalho podem ajudar a validar --- ou a refutar --- previsões de
um artigo em preparação. Alunos cujas contribuições forem incorporadas ao artigo poderão
ser \textbf{convidados como coautores}. Os critérios são os usuais em publicações
científicas:
\begin{enumerate}
\item \textbf{contribuição substancial e original} que sustenta uma afirmação do artigo
--- por exemplo, um experimento, uma figura, um código de análise ou uma explicação para
um desvio da teoria;
\item \textbf{reprodutibilidade}: os professores conseguem regenerar o resultado a
partir do repositório, seguindo o \texttt{README};
\item \textbf{participação na redação e revisão} da parte correspondente e aprovação da
versão final do texto;
\item \textbf{disponibilidade para continuar depois do fim da disciplina}: escrever um
artigo costuma levar meses e exige novas rodadas de experimentos.
\end{enumerate}
Contribuições menores serão reconhecidas nos agradecimentos. A coautoria não depende da
nota, e a nota não depende da coautoria. Os dados de um grupo só serão usados no artigo
com o consentimento de seus autores e com crédito.

\paragraph{Integridade.} Dados fabricados ou “ajustados” para bater com a teoria
resultam em nota zero no trabalho. O uso de ferramentas de IA é permitido para código e
texto, desde que declarado no relatório; o grupo é responsável por tudo o que entrega e
deve ser capaz de explicar cada linha na apresentação.

% =====================================================================
\appendix
\section{Dicas técnicas}
\label{app:dicas}

\paragraph{Amostrar a janela.} O \texttt{cwnd\_monitor.py} roda \texttt{ss -tin} no host
emissor a cada 10\,ms (os hosts do Mininet compartilham o kernel, mas cada um tem seu
\emph{namespace} de rede). Uma alternativa mais precisa é o tracepoint
\texttt{tcp:tcp\_probe} (\path{/sys/kernel/debug/tracing/events/tcp/tcp_probe}), que
registra cwnd e srtt a cada ACK.

\paragraph{Atrasos.} O \texttt{delay} de um \texttt{TCLink} é aplicado em cada sentido;
um enlace com \texttt{delay='10ms'} contribui 20\,ms para o RTT. Use \texttt{params1} e
\texttt{params2} se quiser parâmetros diferentes em cada ponta, por exemplo para colocar
o buffer limitado (\texttt{max\_queue\_size}) só na saída do gargalo.

\paragraph{Offloads.} TSO/GSO/GRO agrupam pacotes e distorcem cwnd e fila. Desligue-os
nas interfaces dos hosts e do gargalo:
\texttt{ethtool -K h1-eth0 tso off gso off gro off}.

\paragraph{Filas.} Inspecione a qdisc real com \texttt{tc -s qdisc show dev s1-eth1}.
Com \texttt{enable\_red=True} o Mininet instala uma qdisc RED com parâmetros fixos;
anote-os no \texttt{metadata.json}.

\paragraph{Detectar perdas.} Uma mudança de \texttt{ssthresh} marca uma redução de
janela (a cwnd em si cai gradualmente, por causa do PRR). Descarte reduções
causadas por \emph{timeout} (cwnd vai a 1 ou 2) e conte-as à parte.

\paragraph{CPU.} Se a vazão medida ficar abaixo de $C$ sem perdas, desconfie da CPU.
Reduza $C$, reduza $n$ ou aumente os recursos da VM.

\section{Código fornecido (\texttt{ig\_starter})}
\label{app:codigo}

\begin{center}
\begin{tabularx}{\textwidth}{@{}lX@{}}
\toprule
Arquivo & O que faz\\
\midrule
\texttt{dumbbell.py} & Topologia haltere com $n$ emissores, RTTs por fluxo, gargalo com
$C$, $B$, RED e perda aleatória opcionais; inicia os fluxos escalonados, os monitores de
cwnd e grava CSVs e \texttt{metadata.json}. \texttt{sudo python3 dumbbell.py -h} lista as
opções.\\
\texttt{cwnd\_monitor.py} & Amostra \texttt{ss -tin} no host emissor e grava
\texttt{t\_s, flow\_id, cc, cwnd\_pkts, ssthresh\_pkts, srtt\_ms, bytes\_acked}.\\
\texttt{ig\_metrics.py} & Injustiça $\mathcal U$, ineficiência $\mathcal I$, $\Dis$,
leitura dos CSVs. Contém dois \texttt{TODO} que são parte do trabalho:
\texttt{detect\_backoffs} (Parte 4) e \texttt{sync\_estimate} (Parte 6).\\
\texttt{plot\_cwnd.py} & Gráfico de cwnd por fluxo (pergunta 4.1 e Parte 2).\\
\texttt{run\_p4.sh} & Exemplo de varredura da Parte 4 com repetições.\\
\bottomrule
\end{tabularx}
\end{center}

Exemplo (Parte 4, $n=2$, Reno, $B=100$, RTT de 40\,ms, fluxos escalonados de 15\,s):
\begin{lstlisting}[style=sh]
sudo python3 dumbbell.py --n 2 --cc reno --rtt-ms 40 --bw 10 --buffer 100 \
     --stagger 15 --duration 150 --out dados/p4/reno-n2-B100/rep1
python3 plot_cwnd.py dados/p4/reno-n2-B100/rep1 -o p4-cwnd-reno-n2-B100.png
\end{lstlisting}
Cuidado numérico: quando $\mathcal U<10^{-8}$, o cálculo de $\log\AM-\overline{\log z}$
perde precisão por cancelamento; não interprete razões de ciclo nesse nível.

\end{document}
